\documentclass[11pt]{article}
\usepackage[margin=1.1in]{geometry}
\usepackage{amsmath,amssymb,amsthm}
\usepackage{booktabs}
\usepackage{enumitem}
\usepackage[round]{natbib}
\setcitestyle{aysep={}}
\usepackage{xcolor}
\definecolor{linknavy}{RGB}{25,55,120}
\usepackage[colorlinks=true,linkcolor=linknavy,citecolor=linknavy,urlcolor=linknavy,
pdftitle={Exact and stable games need not be egalitarian stable},
pdfauthor={Maxwell Flanagan}]{hyperref}

\newtheorem{theorem}{Theorem}[section]
\newtheorem{lemma}[theorem]{Lemma}
\newtheorem{proposition}[theorem]{Proposition}
\newtheorem{corollary}[theorem]{Corollary}
\newtheorem{definition}[theorem]{Definition}
\newtheorem{example}[theorem]{Example}
\newtheorem{remark}[theorem]{Remark}

\newcommand{\R}{\mathbb{R}}
\newcommand{\core}{\operatorname{Core}}
\newcommand{\PES}{\operatorname{PES}}

\title{Exact and stable games need not be egalitarian stable}

\author{Maxwell Flanagan\\[2pt]
\normalsize [Affiliation, City, Country --- to be completed]\\
\normalsize Corresponding author: \texttt{maxwell.flanagan@gmail.com}%
\thanks{A machine-readable certificate supplement with a standalone verifier
accompanies this paper; see the verification section for methodology and the
disclosure of AI assistance.}}

\date{July 2026 --- draft}

\begin{document}
\maketitle

\begin{abstract}
Dietzenbacher (Games and Economic Behavior 127, 2021) asked whether every
exact game and every stable-core game is egalitarian stable, that is, whether
the procedural egalitarian solution always selects a core allocation on these
classes; by known results this holds for exact games with at most four players
and stable games with at most five. We answer both halves of the question in
the negative. For exactness: there is a five-player exact game --- monotone,
superadditive, with grand-coalition worth $4$ --- whose egalitarian claims all
equal $1$, so the claims sum to $5>4$ and the procedural egalitarian solution
leaves the core. A padding construction yields exact non--egalitarian-stable
games for every $n\ge5$, so exactness implies egalitarian stability precisely
up to four players. For stability: there is a six-player exact game with a von
Neumann--Morgenstern stable core that is not egalitarian stable, and the
threshold $n=6$ is again sharp. Both witnesses arise from one device: the
pointwise minimum of finitely many efficient vectors is always an exact game.
The stable witness has a one-dimensional core, and for such segment-core games
we prove that core stability admits a finite certificate --- an arrangement-vertex
enumeration --- which we discharge in exact rational arithmetic. We further
isolate the mechanism (a strict over-cover of the claims) and explain why
random sampling is unlikely to find such games.
\end{abstract}

\medskip
\noindent\textbf{Keywords:} cooperative games; core; exact games;
egalitarian stability; procedural egalitarian solution.

\section{Introduction}\label{sec:intro}

Egalitarianism and core stability interact delicately in
transferable-utility (TU) games. The procedural egalitarian solution
\citep{DBH2017} extends the constrained egalitarian solution of
\citet{DuttaRay89} from convex games to all TU games through an iterative
procedure of intercoalitional equal division; \citet{Dietzenbacher2021}
showed that this solution selects a core allocation exactly on the class of
\emph{egalitarian stable} games, proved that all games with a large core
\citep{Sharkey1982} are egalitarian stable, and---noting that by
\citet{Biswas1999} every exact game \citep{Schmeidler1972} with at most four
players has a large core---posed the following question: ``Whether all exact
games and all stable core games with an arbitrary number of players are
egalitarian stable is an interesting open question for future research.''

This paper answers both halves of the question in the negative, each at the
smallest possible size.

\begin{itemize}[leftmargin=2em]
\item \textbf{A five-player counterexample (Section~\ref{sec:witness}).}
There is an exact, monotone, superadditive five-player game with
$v(N)=4$ whose egalitarian claims are $(1,1,1,1,1)$: every player's claim is
locked at $1$ by the procedure, yet the grand coalition can only pay $4$.
The game is generated by a device of independent interest: the pointwise
minimum of finitely many efficient payoff vectors is always an exact game
(Lemma~\ref{lem:generator}). Two further, structurally different witnesses
found by independent searches are reported alongside.
\item \textbf{A sharp threshold (Section~\ref{sec:family}).} A padding
operation produces from any witness an exact non--egalitarian-stable game
with any number $n\ge5$ of players, with a complete proof for null padding.
Combined with the positive result for $n\le4$, exactness implies egalitarian
stability if and only if $n\le 4$.
\item \textbf{A six-player stable counterexample
(Section~\ref{sec:stable}).} There is an exact six-player game with a stable
core that is not egalitarian stable; since largeness and core stability
coincide for $n\le5$ \citep{Estevez2012} and large-core games are egalitarian
stable, the threshold $n=6$ is sharp. Its core is a one-dimensional segment,
and we prove a finite stability certificate for segment-core games
(Proposition~\ref{prop:cert}): the undominated imputations form a finite union
of polytopes whose vertices lie in an explicit plane arrangement, so
enumerating the arrangement's vertices decides stability. The certificate is
discharged computationally in exact rational arithmetic.
\item \textbf{Mechanism and method (Section~\ref{sec:mechanism}).} A game is
egalitarian stable if and only if its claims vector lies in the core
(Lemma~\ref{lem:esiff}); failure requires a strict over-cover of the claims,
which in all known witnesses over-counts a single player. We also explain
why the question resisted sampling: exactified random games almost never
leave the large-core region, where stability is guaranteed; the witnesses
live on the thin-core frontier that the minimum-of-points generator
produces directly.
\end{itemize}

All claims are certified in exact rational arithmetic; the procedure
implementations were validated against the worked examples of
\citet{Dietzenbacher2021}, and every witness was verified through several
independently written implementations
(Section~\ref{sec:methods}).

%% ===================================================================
%% sections_theory.tex  --- plain LaTeX fragments (no preamble).
%% Assumes the paper-1 preamble: amsmath/amssymb/amsthm, booktabs,
%% enumitem, natbib; theorem-like envs {theorem,lemma,proposition,
%% corollary,definition,example,remark}; macros \R, \core.
%% Player i = index i (1-indexed).  §1 Introduction supplied separately.
%% ===================================================================

\section{Preliminaries and the egalitarian procedure}\label{sec:prelim}

A \emph{transferable-utility (TU) game} is a pair $(N,v)$ with $N=\{1,\dots,n\}$
and $v\colon 2^{N}\to\R$, $v(\emptyset)=0$; we write $v$ for the game and
$x(S)=\sum_{i\in S}x_i$ for $x\in\R^{N}$. The \emph{core} is
\[
  \core(v)=\bigl\{x\in\R^{N}\;:\;x(N)=v(N),\ x(S)\ge v(S)\ \forall S\subseteq N\bigr\},
\]
and $v$ is \emph{balanced} when $\core(v)\neq\emptyset$. The game is
\emph{superadditive} if $v(S\cup T)\ge v(S)+v(T)$ for disjoint $S,T$,
\emph{monotone} if $v(S)\le v(T)$ for $S\subseteq T$, and \emph{convex} if
$v(S\cup T)+v(S\cap T)\ge v(S)+v(T)$ for all $S,T$. Following
\citet{Schmeidler1972}, $v$ is \emph{exact} if for every $S\subseteq N$ there is
a core element attaining $v(S)$, equivalently
\[
  \min_{x\in\core(v)} x(S)=v(S)\qquad\text{for all }S\subseteq N .
\]
The core is \emph{large} \citep{Sharkey1982} if every $x\in\R^{N}$ with
$x(S)\ge v(S)$ for all $S$ dominates a core element from below. One has
$\text{convex}\subseteq\text{large core}\subseteq\text{exact}\subseteq
\text{totally balanced}$, with all inclusions strict for $n\ge5$.

Exactness admits a convenient finite generator, which we use to build and
certify witnesses. It is the finite specialization of Schmeidler's
lower-envelope description of exact games (an exact game is the lower envelope
of its own core); we record it explicitly.

\begin{lemma}[finite generators of exact games]\label{lem:generator}
Let $P\subseteq\R^{N}$ be a nonempty finite set of vectors sharing a common sum
$\tau$, and define $v(S)=\min_{x\in P}x(S)$. Then $v(N)=\tau$, every $x\in P$
lies in $\core(v)$, and $v$ is exact.
\end{lemma}

\begin{proof}
Every $x\in P$ has $x(N)=\tau=v(N)$ and $x(S)\ge\min_{y\in P}y(S)=v(S)$ for all
$S$, so $x\in\core(v)$; in particular $\core(v)\neq\emptyset$. Fix $S$ and pick
$x^{\ast}\in\arg\min_{x\in P}x(S)$. Then $x^{\ast}\in\core(v)$ and
$x^{\ast}(S)=v(S)$, so $\min_{x\in\core(v)}x(S)\le v(S)$; the reverse inequality
holds because every core element satisfies $x(S)\ge v(S)$. Hence
$\min_{\core(v)}x(S)=v(S)$ for all $S$.
\end{proof}

\paragraph{The egalitarian procedure.}
We recall the procedure underlying the \emph{procedural egalitarian solution}
\citep{DBH2017,Dietzenbacher2021}. Set $P^{v,0}=\emptyset$. For $k\ge1$ and
$S\neq\emptyset$ define $\chi^{v,k}(S)\in\R^{S}$ by
\[
  \chi^{v,k}_i(S)=
  \begin{cases}
    \gamma^{v,k-1}_i, & i\in S\cap P^{v,k-1},\\[4pt]
    \dfrac{v(S)-\sum_{j\in S\cap P^{v,k-1}}\gamma^{v,k-1}_j}
          {\lvert S\setminus P^{v,k-1}\rvert}, & i\in S\setminus P^{v,k-1},
  \end{cases}
\]
i.e.\ already-fixed members keep their payoff and the residual worth of $S$ is
split equally among the rest. The \emph{egalitarian admissible} collection of
round $k$ is
\[
  A^{v,k}=\Bigl\{S\;:\;\textstyle\sum_{i\in S}\chi^{v,k}_i(S)=v(S)\ \text{ and }\
  \forall i\in S\ \forall T\ni i:\ \chi^{v,k}_i(T)\le\chi^{v,k}_i(S)\Bigr\},
\]
so $S$ is admissible when it is efficient for $\chi^{v,k}$ and gives each of its
members their pointwise-largest $\chi^{v,k}$-value. Put
$P^{v,k}=\bigcup_{S\in A^{v,k}}S$ and $\gamma^{v,k}_i=\chi^{v,k}_i(S)$ for
$i\in P^{v,k}$, $S\in A^{v,k}\ni i$. The procedure terminates at
$n_v=\min\{k:P^{v,k}=N\}$; the \emph{egalitarian claims} vector is
$\widehat\gamma^{\,v}=\gamma^{v,n_v}$ and $\widehat A^{v}=A^{v,n_v}$.

The claims vector is always an \emph{aspiration} \citep{Bennett1983}:
$\widehat\gamma^{\,v}(S)\ge v(S)$ for all $S$, and for each $i\in N$ there is
$S\ni i$ with $\widehat\gamma^{\,v}(S)=v(S)$ and
$\widehat\gamma^{\,v}_i\le\widehat\gamma^{\,v}_j$ for all $j\in S$. A game is
\emph{egalitarian stable} (ES) if $N\in\widehat A^{v}$; equivalently, the
procedural egalitarian solution then selects from the core. Convex and
large-core games are ES \citep{Dietzenbacher2021}. Since exact games with
$n\le4$ have large cores \citep{Biswas1999}, they are ES; whether every exact
game is ES was left open \citep[][Sec.~4]{Dietzenbacher2021}. We resolve it
negatively.


\section{ES equals membership of the claims vector in the core}\label{sec:lemma1}

The following characterization reduces egalitarian stability of any balanced
game to a single membership test and organizes everything below. Write
$g(S)=\widehat\gamma^{\,v}(S)-v(S)\ge 0$ for the aspiration gap. A collection of
weights $\lambda=(\lambda_S)_{S\subsetneq N}$, $\lambda_S\ge0$, is a
\emph{fractional cover} if $\sum_{S\ni i}\lambda_S\ge1$ for every $i\in N$.

\begin{lemma}\label{lem:esiff}
Let $v$ be balanced with egalitarian claims $\widehat\gamma=\widehat\gamma^{\,v}$.
The following are equivalent:
\begin{enumerate}[label=\textup{(\roman*)},leftmargin=2.4em]
\item $v$ is egalitarian stable;
\item $\widehat\gamma\in\core(v)$;
\item $\sum_{i\in N}\widehat\gamma_i=v(N)$;
\item $\widehat\gamma$ dominates a core element from below;
\item $g(N)\le\min\bigl\{\sum_{S\subsetneq N}\lambda_S\,g(S)\;:\;\lambda
      \text{ a fractional cover}\bigr\}$.
\end{enumerate}
\end{lemma}

\begin{proof}
$\widehat\gamma$ is an aspiration, so $\widehat\gamma(S)\ge v(S)$ for all $S$ and
$\sum_{i\in N}\widehat\gamma_i\ge v(N)$ always.

(i)$\Rightarrow$(iii). If $v$ is ES then $N\in\widehat A^{v}$, and the efficiency
requirement in the definition of $A^{v,n_v}$ gives
$\sum_{i\in N}\chi^{v,n_v}_i(N)=v(N)$; at the terminal round every player is
fixed, so $\chi^{v,n_v}_i(N)=\widehat\gamma_i$ and $\sum_i\widehat\gamma_i=v(N)$.

(iii)$\Rightarrow$(ii). With $\sum_i\widehat\gamma_i=v(N)$ and
$\widehat\gamma(S)\ge v(S)$ for all $S$, $\widehat\gamma$ is efficient and
satisfies every core inequality, i.e.\ $\widehat\gamma\in\core(v)$.

(ii)$\Rightarrow$(iv) is trivial. (iv)$\Rightarrow$(i): let $y\in\core(v)$ with
$y\le\widehat\gamma$. For every $S\in\widehat A^{v}$,
$y(S)\le\widehat\gamma(S)=v(S)\le y(S)$, so $y(S)=\widehat\gamma(S)$ and, as
$y\le\widehat\gamma$, $y_i=\widehat\gamma_i$ for $i\in S$. The admissible
coalitions cover $N$, so $y=\widehat\gamma\in\core(v)$; at the terminal round
$N$ is then admissible, i.e.\ $v$ is ES.

(iv)$\Leftrightarrow$(v). With $y=\widehat\gamma-t$, domination of a core element
is feasibility of $t\ge0$, $t(N)=g(N)$, $t(S)\le g(S)$ ($S\subsetneq N$); this
holds iff $\max\{t(N):t\ge0,\ t(S)\le g(S)\}\ge g(N)$, and LP duality identifies
the maximum with $\min\{\sum_S\lambda_S g(S):\lambda\text{ a fractional cover}\}$.
\end{proof}

\begin{remark}\label{rem:overcover}
By Bondareva--Shapley, every fractional \emph{partition} (one with
$\sum_{S\ni i}\lambda_S=1$) obeys $\sum_S\lambda_S v(S)\le v(N)$, and hence
$\sum_S\lambda_S g(S)\ge g(N)$. A failure of \textup{(v)} must therefore be
certified by a \emph{strict over-cover}, one in which some player is covered with
total weight exceeding~$1$. This is the analytic signature of the witnesses
below: the egalitarian claims inflate precisely because one pivotal player is
over-covered by the tight coalitions that fix the claims.
\end{remark}


\section{A five-player exact game that is not egalitarian stable}
\label{sec:witness}

\begin{theorem}\label{thm:witness}
There is an exact five-player game that is not egalitarian stable. Consequently
the implication ``exact $\Rightarrow$ ES,'' valid for $n\le4$, fails at $n=5$.
\end{theorem}

\begin{proof}
Let $N=\{1,2,3,4,5\}$ and take the two allocations
\[
  x^{1}=(0,0,1,1,2),\qquad x^{2}=(1,1,0,2,0),
\]
both with total $4$, and set $v(S)=\min\{x^{1}(S),x^{2}(S)\}$ for all $S$. By
Lemma~\ref{lem:generator}, $v$ is exact, with $x^{1},x^{2}\in\core(v)$ and
$v(N)=4$; the game is monotone and superadditive but not convex, and its full
coalition function is listed in the supplement. In fact $x^{1}$ and $x^{2}$ are
the \emph{only} vertices of $\core(v)$, so the core is the one-dimensional
segment $\operatorname{conv}\{x^{1},x^{2}\}$---the thinnest nondegenerate core,
and exactly the regime in which non-egalitarian-stability concentrates
(Remark~\ref{rem:frequency}).

\emph{The procedure.} The maximum average worth is
$\max_{S}v(S)/\lvert S\rvert=1$, attained exactly on
\[
  \widehat A^{v}=\bigl\{\{4\},\{3,4\},\{4,5\},\{1,4,5\},\{2,4,5\}\bigr\},
\]
each of which routes through the pivot player $4$ (see the single-round trace,
Table~\ref{tab:trace}). Every player attains its own maximum average, $1$, on one
of these coalitions, so all five are egalitarian admissible in the first round;
thus $n_v=1$ and
\[
  \widehat\gamma^{\,v}=(1,1,1,1,1),\qquad \sum_{i\in N}\widehat\gamma_i=5>4=v(N).
\]
By Lemma~\ref{lem:esiff} the game is not egalitarian stable: the grand coalition
would give each player only $v(N)/5=\tfrac45<1$, so $N\notin\widehat A^{v}$.
\end{proof}

\begin{table}[t]
\centering
\caption{Egalitarian procedure for the primary witness (single round, cf.\
\citealp{Dietzenbacher2021}, Ex.~2). Entries are $\chi^{v,1}_i(S)$ over
$i=1,\dots,5$, ``$\cdot$'' for $i\notin S$; the last column is $N$.}
\label{tab:trace}
\begin{tabular}{lcccccc}
\toprule
$S$ & $\{4\}$ & $\{3,4\}$ & $\{4,5\}$ & $\{1,4,5\}$ & $\{2,4,5\}$ & $N$\\
$v(S)$ & $1$ & $2$ & $2$ & $3$ & $3$ & $4$\\
\midrule
$\chi^{v,1}(S)$
 & $(\cdot,\cdot,\cdot,1,\cdot)$
 & $(\cdot,\cdot,1,1,\cdot)$
 & $(\cdot,\cdot,\cdot,1,1)$
 & $(1,\cdot,\cdot,1,1)$
 & $(\cdot,1,\cdot,1,1)$
 & $(\tfrac45,\tfrac45,\tfrac45,\tfrac45,\tfrac45)$\\
\bottomrule
\end{tabular}
\end{table}

Here the egalitarian claims are \emph{symmetric}, yet the game is not ES: what
inflates them is not an unequal claim but the fact that every tight coalition is
routed through the single pivot $4$ (the over-cover of
Remark~\ref{rem:overcover}). The inclusion-maximal admissible coalitions
$\{3,4\},\{1,4,5\},\{2,4,5\}$ have $\{4\}$ as their only common member, so the
strong egalitarian claimants are $D^{v}=\{4\}$ and the procedural egalitarian
solution equalizes the rest down to $\lambda=\tfrac34$,
\[
  PES(v)=\bigl(\tfrac34,\tfrac34,\tfrac34,1,\tfrac34\bigr),
\]
which leaves the core, violating $\{3,4\}$ and $\{4,5\}$ by $\tfrac14$ and
$\{1,4,5\}$ and $\{2,4,5\}$ by $\tfrac12$. A core allocation attaining $v(S)$ for
each of the $30$ proper coalitions (the exactness certificate) is provided in the
supplement.

\begin{remark}[a minimal-mass variant]\label{rem:minmass}
The grand-coalition worth $v(N)=4$ is minimal over two-point generators, and the
five tight coalitions that drive the procedure keep their worths under
coalitionwise reduction. Greedily lowering the non-tight worths yields a variant of
total worth-mass $41$ (from $46$) that is again exact, monotone and superadditive,
has the identical procedure trace (Table~\ref{tab:trace}), the same claims, and the
same $\PES$ with the same four core violations, and is locally minimal---no single
worth admits a further decrease preserving these properties (supplement). Its core
has seven vertices, so the two-point generator and the segment core are particular
to the reference game above.
\end{remark}

\begin{remark}\label{rem:notlarge}
By Lemma~\ref{lem:esiff} the witness is exact but \emph{not} large-core, so it
lies in the gap between exact and large-core games---empty for $n\le4$, first
opening at $n=5$ \citep{Biswas1999}. Two further, independently located witnesses confirm
the phenomenon and its two modes: an exact game with worths at most $6$ and an
\emph{asymmetric} pivot claim, $\widehat\gamma=(1,1,3,1,1)$, where pivot player
$3$ carries claim $3$ and $\sum_i\widehat\gamma_i=7>6=v(N)$; and an exact game
whose claims $\tfrac12(2,7,3,7,2)$ sum to $\tfrac{21}{2}>10=v(N)$. All three were
mutually cross-verified across independent implementations; their data are in the
supplement.
\end{remark}


\section{An infinite family: exact non-ES games for every $n\ge5$}
\label{sec:family}

We upgrade Theorem~\ref{thm:witness} to a family covering all $n\ge5$ by a
padding operation that provably preserves both exactness and non-ES.

\begin{definition}\label{def:pad}
Let $v$ be a game on $N$ and $c\ge0$. The \emph{$c$-padding} of $v$ is the game
$v^{+}$ on $N\cup\{p\}$ ($p\notin N$) with $v^{+}(S)=v(S)$ for $S\subseteq N$ and
$v^{+}(S\cup\{p\})=v(S)+c$. For $c=0$, $p$ is a \emph{null} player; for $c\ge1$,
$p$ is essential.
\end{definition}

\begin{lemma}[padding preserves exactness]\label{lem:padexact}
If $v$ is exact then so is every $c$-padding $v^{+}$ ($c\ge0$).
\end{lemma}

\begin{proof}
Let $x\in\core(v^{+})$. From $x(N)\ge v(N)$, $x_p\ge v^{+}(\{p\})=c$ and
$x(N)+x_p=v^{+}(N)=v(N)+c$ we get $x(N)=v(N)$ and $x_p=c$; conversely any
$(y,c)$ with $y\in\core(v)$ lies in $\core(v^{+})$. Hence
$\core(v^{+})=\core(v)\times\{c\}$, and
$\min_{\core(v^{+})}x(S)=\min_{\core(v)}y(S\setminus\{p\})+c\,[p\in S]
=v(S\setminus\{p\})+c\,[p\in S]=v^{+}(S)$ for every $S$.
\end{proof}

\begin{lemma}[null players are inert]\label{lem:inert}
Let $v$ be nonnegative and $v^{+}$ its $0$-padding by a null player $p$. Then the
egalitarian procedure on $v^{+}$ fixes each $i\in N$ in the same round and at the
same value as on $v$, and fixes $p$ last at value $0$; in particular
$\widehat\gamma^{\,v^{+}}=(\widehat\gamma^{\,v},0)$.
\end{lemma}

\begin{proof}
By induction on $k$: $P^{v^{+},k}\setminus\{p\}=P^{v,k}$, $\gamma^{v^{+},k}$
agrees with $\gamma^{v,k}$ on $P^{v,k}$, and $p\notin P^{v^{+},k}$ while
$P^{v,k}\neq N$. The base is trivial. Assume it for $k-1$, with $Q=P^{v,k-1}$,
claims $\gamma$, and $p$ free. For $S\subseteq N$ the data defining
$\chi^{v^{+},k}(S)$ and $\chi^{v,k}(S)$ coincide, so they are equal. For
$R=S\cup\{p\}$ with $f=\gamma(Q\cap S)$: if $S\setminus Q\neq\emptyset$, each free
member of $R$ gets $\sigma=(v(S)-f)/(|S\setminus Q|+1)$, whereas
$\chi^{v^{+},k}_i(S)=(v(S)-f)/|S\setminus Q|$ for $i\in S\setminus Q$; if
$v(S)-f\ge0$ then $\sigma\le\chi^{v^{+},k}_i(S)$, and if $v(S)-f<0$ then
$\sigma<0\le v(\{i\})\le\max_{T\ni i}\chi^{v^{+},k}_i(T)$ (nonnegativity, with
$\{i\}$ available). Either way no original free member attains a new maximum at
$R$, so no such $R$ is admissible in round $k$; the only admissible coalitions
containing $p$ have all original members already in $Q$, and they fix $p$ alone.
Hence the newly fixed original players match those of $v$, proving the step. At
termination $\widehat\gamma^{\,v^{+}}_p=\max_{S\subseteq
N}(v(S)-\widehat\gamma^{\,v}(S))\le0$ (aspiration) and $\ge v^{+}(\{p\})=0$, so
$=0$.
\end{proof}

\begin{theorem}\label{thm:family}
For every $n\ge5$ there is an exact game on $n$ players that is not egalitarian
stable. Concretely, padding the primary witness of Theorem~\ref{thm:witness}
with $n-5$ null players yields such a game.
\end{theorem}

\begin{proof}
The witness $v$ ($n=5$) is exact, nonnegative, and non-ES with
$\sum_i\widehat\gamma^{\,v}_i=5>4=v(N)$. Adjoin a null player: by
Lemma~\ref{lem:padexact} the padding $v^{+}$ is exact, and by
Lemma~\ref{lem:inert} $\widehat\gamma^{\,v^{+}}=(\widehat\gamma^{\,v},0)$, so
$\sum_i\widehat\gamma^{\,v^{+}}_i=5>4=v^{+}(N)$; by Lemma~\ref{lem:esiff}
$v^{+}$ is not ES. Nonnegativity and both properties are preserved, so iterating
$n-5$ times gives an exact non-ES game for each $n\ge5$.
\end{proof}

Null players are convenient for the proof but degenerate. Padding with
\emph{essential} players ($c\ge1$) gives non-degenerate examples with the same
inflation; we verified (supplement) that for $c=1$ the padded player receives
claim $c$, all other claims are unchanged, and the games remain exact and non-ES
through $n=9$, with every player essential.

\begin{corollary}[sharp threshold]\label{cor:threshold}
Every exact game with $n\le4$ is egalitarian stable, and for every $n\ge5$ there
is an exact game that is not. The threshold $n=5$ coincides with the onset of the
exact/large-core gap.
\end{corollary}

\begin{proof}
For $n\le4$, exact games have large cores \citep{Biswas1999} and large-core games
are ES \citep{Dietzenbacher2021}. For $n\ge5$, Theorem~\ref{thm:family}.
\end{proof}


\section{The mechanism, and why the question stayed open}\label{sec:mechanism}

Lemma~\ref{lem:esiff} and Remark~\ref{rem:overcover} isolate the combinatorial
cause of non-ES. We record it as a sufficient condition covering both observed
modes, then explain why a naive symmetric family cannot realize it and why the
phenomenon eluded random sampling.

\begin{proposition}[pivot over-cover]\label{prop:hub}
Let $v$ be exact, and suppose there is a player $h$ and tight coalitions
$T_1,\dots,T_r\ni h$ (each $\widehat\gamma(T_j)=v(T_j)$) that cover $N$, are
each fixed by the egalitarian procedure, and each spend $h$'s full claim
$\widehat\gamma_h$. Then $\{T_j\}$ is a strict over-cover of $N$ at $h$, whence
$\sum_i\widehat\gamma_i>v(N)$ and $v$ is not egalitarian stable.
\end{proposition}

\begin{proof}
Each $T_j$ is tight, so unit weights on the $T_j$ give a fractional cover with
$\sum_j g(T_j)=0$; as every $T_j$ contains $h$, the cover assigns $h$ total
weight $\ge2$, i.e.\ it is strict, and Lemma~\ref{lem:esiff}\,(v) fails once
$g(N)>0$. Equivalently, summing $\widehat\gamma(T_j)=v(T_j)$ double-counts
$\widehat\gamma_h$, forcing $\sum_i\widehat\gamma_i>v(N)$.
\end{proof}

Proposition~\ref{prop:hub} unifies the two modes seen in
Remark~\ref{rem:notlarge}. In the \emph{asymmetric} mode (the hub witness) the
pivot $h=3$ additionally carries an inflated claim $\widehat\gamma_h=3$. In the
\emph{symmetric} mode (the primary witness) all claims equal $1$, yet the five
tight coalitions still route through the single pivot $h=4$; the symmetric mode
is exactly the special case of equal claims. In both, exactness survives because
``absorber'' players let the core attain the pivot's small-coalition worths (for
the primary witness, $x^{1}$ and $x^{2}$ attain the relevant worths from opposite
sides): the tight coalitions certifying the claims are pairwise
incompatible---all require the pivot---yet each individual worth remains
core-attainable.

\begin{remark}[symmetric scaling collapses to ES]\label{rem:obstruction}
A tempting symmetric scaling---a hub $h$ with $v(\{h\})=m$, one absorber $D$, and
$k$ leaves, with $v(S)=m+2$ for $h,D\in S$ and $\ge1$ leaf, $v(\{h,D\})=m+1$,
$v(N)=m+3$---is non-ES for all $k\ge3$ (claims $(m,1,\dots,1)$, sum $m+1+k>m+3$)
but is \emph{not exact}: $\min_{\core}x(\{h,D\})=m+2>m+1$, since
$x(\{h,D,\text{leaf}\})\ge m+2$ with $x(N)=m+3$ pins every leaf to $0$. Raising
$v(\{h,D\})$ to its core-minimum---the exactification forced by exactness---hands
the surplus to the absorber (claim $\tfrac32$, leaves $\tfrac12$) and the game
becomes ES; its $n=4$ member is exact \emph{and} ES, as
Corollary~\ref{cor:threshold} requires. A genuine exact non-ES game must
therefore be asymmetric---several satellites and absorbers, as in the
witnesses---which is why padding an existing witness, not scaling a symmetric
template, is the robust route to Theorem~\ref{thm:family}.
\end{remark}

\begin{remark}[why the question stayed open]\label{rem:frequency}
Non-ES exact games are not pathological once one looks in the right place, but
they are nearly invisible to the natural sampling scheme. Exactifying random
balanced games essentially never leaves the large-core class (in our sampling,
$3$ of $499$ exactified games were even non-large-core, and a random-aspiration
probe detected non-largeness in none), so such a search almost never produces a
non-ES game. By contrast, drawing exact games as lower envelopes of a few points
(Lemma~\ref{lem:generator}) lands on the thin-core frontier where the phenomenon
concentrates: among $18{,}507$ point-generated exact five-player games, $493$
($2.66\%$) were not egalitarian stable. Within these, the symmetric-claims mode of
the reference witness is itself rare---only $2$ of the $493$ have all claims
equal---and the claim inflation $\sum_i\widehat\gamma_i-v(N)$ concentrates at the
values $1$ and $\tfrac12$. The offending aspiration is always the egalitarian
claims vector itself, which generic aspiration sampling misses; this mismatch
between the natural sampling scheme and the location of the counterexamples is, we
suspect, why the question remained open.
\end{remark}

\section{Stable games need not be egalitarian stable}\label{sec:stable}

We now resolve the second half of the question. Recall that an
\emph{imputation} is an efficient, individually rational allocation, and that an
allocation $x$ \emph{dominates} an imputation $y$ via $S$ if $x_i>y_i$ for all
$i\in S$ and $x(S)\le v(S)$ \citep{vNM1944}. The core, always internally
stable, is a \emph{(von Neumann--Morgenstern) stable set} precisely when it is
also externally stable: every imputation outside the core is dominated by a
core element. We then say the game has a \emph{stable core}. Large-core games
have stable cores \citep{Sharkey1982}, and for games with at most five players
largeness and core stability are equivalent \citep[Sec.~5]{Estevez2012}; since
large-core games are egalitarian stable \citep{Dietzenbacher2021}, every stable
game with $n\le5$ is egalitarian stable. \citet{Dietzenbacher2021} asked
whether this persists for all $n$. It does not:

\begin{theorem}\label{thm:stable}
There is a six-player exact game with a stable core that is not egalitarian
stable. Hence core stability does not imply egalitarian stability, and by the
preceding paragraph the threshold $n=6$ is sharp.
\end{theorem}

The witness is again a two-point generator game. Let $N=\{1,\dots,6\}$,
\[
  A=(1,1,0,1,3,0),\qquad B=(2,0,1,0,1,2),
\]
both with total $6$, and $w_1(S)=\min\{A(S),B(S)\}$; the full coalition
function is in the supplement. A second witness
$w_2=\min\{(3,0,2,2,1,0),(1,1,1,1,2,2)\}$ (total $8$), found independently,
confirms the phenomenon; all statements below are certified for both.

\begin{proof}[Proof of Theorem~\ref{thm:stable}, non-ES half]
By Lemma~\ref{lem:generator}, $w_1$ is exact with $A,B\in\core(w_1)$ and
$w_1(N)=6$; it is monotone. The procedure terminates in one round with
\[
  \widehat A^{w_1}=\bigl\{\{1,2\},\{1,4\},\{1,5\},\{5,6\},\{1,3,5\},
  \{1,2,5,6\},\{1,4,5,6\}\bigr\},
\]
egalitarian claims
$\widehat\gamma=\bigl(\tfrac32,\tfrac12,1,\tfrac12,\tfrac32,\tfrac32\bigr)$,
and $\sum_i\widehat\gamma_i=\tfrac{13}2>6=w_1(N)$; by Lemma~\ref{lem:esiff},
$w_1$ is not egalitarian stable. (For $w_2$: claims
$\bigl(\tfrac32,1,\tfrac32,\tfrac32,\tfrac32,\tfrac32\bigr)$, sum
$\tfrac{17}2>8$.)
\end{proof}

The stability half is where the new difficulty lies: external stability is a
statement about \emph{every} imputation outside the core, and no finite
certificate for it exists in general. The witnesses, however, have the thinnest
possible nondegenerate cores, and for such games we prove that stability
\emph{does} admit a finite certificate.

\paragraph{Segment cores.} Exactly $12$ coalitions $S$ satisfy
$A(S)=B(S)=w_1(S)$ (we call these \emph{fully tight}); together with
efficiency they force affine dimension $1$, so
$\core(w_1)=\operatorname{conv}\{A,B\}=\{x_t:=(1-t)A+tB\;:\;t\in[0,1]\}$, a
segment (certificate in the supplement; likewise for $w_2$). Since
$w_1=\min\{A,B\}$ pointwise, every coalition is tight at $A$ or at $B$.

\begin{lemma}[domination from a segment]\label{lem:segdom}
Let the core be the segment $[A,B]$ and let $S\subsetneq N$. The core elements
attaining $v(S)$ are: the whole segment if $S$ is fully tight; the endpoint
$A$ if $A(S)=v(S)<B(S)$; the endpoint $B$ if $B(S)=v(S)<A(S)$. Consequently an
imputation $y$ is dominated via $S$ only by a tight endpoint or, for fully
tight $S$, by an interior point $x_t$.
\end{lemma}

\begin{proof}
$x_t(S)$ is affine in $t$ and $\ge v(S)$ on $[0,1]$, so its minimum over the
segment is attained at an endpoint; it equals $v(S)$ exactly as listed. A
dominating core element via $S$ must satisfy $x(S)=v(S)$ (as $x(S)\le v(S)$
and $x\in\core$), hence lies in the listed set.
\end{proof}

\begin{lemma}[the domination boundary is cut by finitely many planes]
\label{lem:planes}
For $S\subsetneq N$ define
$\operatorname{opt}_S(y)=\max_{x\in F_S}\min_{i\in S}(x_i-y_i)$, where
$F_S$ is the set of core elements attaining $v(S)$, so that $y$ is dominated
via $S$ iff $\operatorname{opt}_S(y)>0$. Then $\operatorname{opt}_S$ is a
concave, piecewise-affine function of $y$ whose affine pieces have boundaries
contained in the planes
\[
  y_i=A_i,\qquad y_i=B_i,\qquad
  (y_i-A_i)(B_j-A_j)=(y_j-A_j)(B_i-A_i)\quad(i\neq j).
\]
\end{lemma}

\begin{proof}
$\min_{i\in S}(x_i-y_i)$ is a minimum of affine functions of $(t,y)$ and $F_S$
is a segment or a point (Lemma~\ref{lem:segdom}), so $\operatorname{opt}_S$ is
concave and piecewise affine. If $F_S$ is an endpoint, the pieces of
$\min_{i\in S}(x_i-y_i)$ change where two arguments tie or where the minimizing
coordinate changes sign, i.e.\ across planes $y_i=A_i$ or $y_i=B_i$. If $F_S$
is the whole segment, the interior maximizer over $t$ is characterized by two
coordinates $i,j\in S$ tying, $x_{t,i}-y_i=x_{t,j}-y_j$; eliminating $t$ turns
both the optimality boundary and the zero locus of the resulting affine piece
into the stated pairwise plane for $(i,j)$. A tie of three coordinates lies in
the intersection of two pairwise planes, so pairwise planes suffice.
\end{proof}

\begin{proposition}[finite stability certificate for segment cores]
\label{prop:cert}
Let $\mathcal P$ consist of the efficiency plane, the coalition planes
$y(S)=v(S)$, the floor planes $y_i=v(\{i\})$, and the planes of
Lemma~\ref{lem:planes}. If no vertex of the arrangement of $\mathcal P$ within
the imputation set is an undominated imputation outside the core, then the core
is stable.
\end{proposition}

\begin{proof}
Let $U$ be the set of undominated imputations, so $\core(v)\subseteq U$ and
stability is $U=\core(v)$. By Lemma~\ref{lem:planes},
$\{y:\operatorname{opt}_S(y)\le0\}$ is a sublevel set of a concave
piecewise-affine function, i.e.\ a finite union of polyhedra delimited by
planes of $\mathcal P$; intersecting over $S$ and with the imputation set, $U$
is a finite union of polytopes, each with all facets in $\mathcal P$. Suppose
$U\supsetneq\core(v)$ and pick a polytope piece $P_c\subseteq U$ with
$P_c\not\subseteq\core(v)$. The core is convex and cut out by inequalities
$y(S)\ge v(S)$, so if every vertex of $P_c$ lay in the core, then
$P_c=\operatorname{conv}(\text{vertices})\subseteq\core(v)$; hence some vertex
$w$ of $P_c$ --- a vertex of the $\mathcal P$-arrangement --- satisfies
$w\in U\setminus\core(v)$.
\end{proof}

\begin{proof}[Proof of Theorem~\ref{thm:stable}, stability half
(computer-assisted)]
For $w_1$ the family $\mathcal P$ has $89$ planes. We enumerated all
$\binom{89}{5}=41{,}507{,}642$ potential vertex bases (with the efficiency
plane always active); about $20.8$ million are nonsingular, of which about
$4.5$ million solve to imputations outside the core. For \emph{every} such
vertex an exact rational dominating certificate was produced: a core element
$x^{*}$ with $x^{*}(S)=v(S)$ and $x^{*}_i>y_i$ for all $i\in S$, verified in
exact arithmetic. No undominated vertex exists, so by
Proposition~\ref{prop:cert} the core of $w_1$ is stable; the run for $w_2$ is
identical in outcome. Two soundness guarantees accompany the computation.
First, a vertex is discarded as dominated only on an exact rational witness as
above; vertices with small floating-point domination margins are passed to a
full exact undomination check, so no false negatives are possible. Second, the
float classification of a vertex as an imputation outside the core is itself
exact on this vertex set: by Cramer's rule every vertex coordinate is a
rational with denominator bounded by the maximal basis determinant
($\le7\,000$ here by Hadamard's inequality), so every nonzero constraint value
has magnitude at least $1.4\times10^{-4}$, far exceeding both the
classification tolerance ($10^{-7}$) and the solve error ($\approx10^{-12}$).
The enumeration architecture was gated before use: it re-finds the known
escapes of two non-stable segment-core control games (one five-player, one
six-player), including an escape lying on the new plane classes. Full run
statistics, checkpoints, and code are in the supplement.
\end{proof}

\begin{corollary}[full resolution]\label{cor:full}
Neither exactness nor core stability implies egalitarian stability. The
thresholds are sharp: exact games are egalitarian stable precisely up to
$n=4$ (Corollary~\ref{cor:threshold}), and stable games precisely up to
$n=5$.
\end{corollary}

\begin{remark}[why the stable half is harder, and where the witnesses live]
\label{rem:stablehard}
Non-ES requires inflated claims, which by Remark~\ref{rem:frequency}
concentrate on thin cores; but thin cores are generically \emph{not} stable
--- in our experiments, escape imputations were certified for every one of
$146$ randomly generated non-ES games at $n\in\{5,6\}$. The witnesses sit at
the extreme point of this trade-off: a one-dimensional core that nevertheless
dominates every outside imputation. This coupling --- thinness driving
non-ES while destroying stability --- is also why the stable half could not be
settled by sampling, and why the witnesses had to be certified rather than
searched. Both witnesses are exact but not large-core (automatic, since
large-core games are ES), so they simultaneously exhibit the exact/ES and the
stable/large gaps at $n=6$.
\end{remark}

%% Verification note (fold into the assembled methods section): all quantitative
%% claims are certified in exact rational arithmetic. The machine-readable
%% supplement (supplement_data.json) records, for the primary witness, the
%% coalition function and its two generator points, a core allocation attaining
%% v(S) for each of the 30 proper coalitions (exactness certificate), the
%% procedure trace, the egalitarian claims, admissible collection, strong
%% claimants, PES and its four core violations; and the analogous data for the two
%% secondary witnesses. The egalitarian procedure was implemented independently by
%% three authors, validated on Examples 2 and 3 of Dietzenbacher (2021), and the
%% three witnesses mutually cross-verified; the family and inertness statements
%% were checked through n=9.


\section{Verification methodology}\label{sec:methods}

All numerical claims are certified in exact rational arithmetic. The
egalitarian procedure was implemented four times, independently, by
different authorship processes; each implementation was gated on exact
reproduction of Examples~2 and~3 of \citet{Dietzenbacher2021} (claims,
admissible collections, strong claimants, and solution values) before any
search was run, and one implementation was additionally calibrated on $235$
random exact four-player games, all of which are egalitarian stable as
required by the results of \citet{Biswas1999} and \citet{Dietzenbacher2021}.
The three witnesses reported here were each discovered independently and
cross-verified through all implementations; exactness is certified by
exhibiting, for every proper coalition, a core allocation attaining its
worth (linear programs solved in exact rational arithmetic by a solver
fuzz-tested against a standard floating-point implementation). The stability certificate of Theorem~\ref{thm:stable} was discharged by
exhaustive enumeration of the $\binom{89}{5}$ vertex bases of the certifying
arrangement for each six-player witness; every outside-core vertex was
discarded only against an exact rational dominating core element, borderline
cases were passed to a full exact undomination check, and the enumeration
architecture was validated by re-finding the known escapes of two non-stable
segment-core control games before use. The six-player witnesses' exactness,
non--egalitarian-stability, and segment-core geometry were additionally
re-derived by a second, independently written toolchain (separate exact LP
solver and separate implementation of the egalitarian procedure). A
machine-readable supplement contains the full certificate data ---
coalition functions, procedure traces, claims, admissible collections,
solution values, core violations, and per-coalition exactness
witnesses --- together with a self-contained verifier using only the Python
standard library. The computational search, the verification tooling, and
the drafting of this paper made substantial use of large-language-model
assistance; every mathematical claim was verified by the certification
pipeline described here, and all text was reviewed by the author.

\subsection*{Statements and Declarations}
\small
\textbf{Funding:} none.
\textbf{Data and code availability:} certificate supplement and verifier
accompany the paper.
\textbf{Competing interests:} none.
\normalsize

\begin{thebibliography}{99}

\bibitem[Bennett(1983)]{Bennett1983}
Bennett E (1983) The aspiration approach to predicting coalition formation
and payoff distribution in sidepayment games.
Int J Game Theory 12:1--28.
\url{https://doi.org/10.1007/BF01756101}

\bibitem[Biswas et al.(1999)]{Biswas1999}
Biswas AK, Parthasarathy T, Potters JAM, Voorneveld M (1999) Large cores
and exactness. Games Econ Behav 28:1--12.
\url{https://doi.org/10.1006/game.1998.0686}

\bibitem[Dietzenbacher(2021)]{Dietzenbacher2021}
Dietzenbacher B (2021) Monotonicity and egalitarianism.
Games Econ Behav 127:194--205.
\url{https://doi.org/10.1016/j.geb.2021.03.006}

\bibitem[Dietzenbacher et al.(2017)]{DBH2017}
Dietzenbacher B, Borm P, Hendrickx R (2017) The procedural egalitarian
solution. Games Econ Behav 106:179--187.
\url{https://doi.org/10.1016/j.geb.2017.10.001}

\bibitem[Dutta and Ray(1989)]{DuttaRay89}
Dutta B, Ray D (1989) A concept of egalitarianism under participation
constraints. Econometrica 57:615--635.
\url{https://doi.org/10.2307/1911055}

\bibitem[Est\'evez-Fern\'andez(2012)]{Estevez2012}
Est\'evez-Fern\'andez A (2012) New characterizations for largeness of the
core. Games Econ Behav 76:160--180.
\url{https://doi.org/10.1016/j.geb.2012.04.003}

\bibitem[von Neumann and Morgenstern(1944)]{vNM1944}
von Neumann J, Morgenstern O (1944) Theory of Games and Economic Behavior.
Princeton University Press, Princeton

\bibitem[Schmeidler(1972)]{Schmeidler1972}
Schmeidler D (1972) Cores of exact games, I.
J Math Anal Appl 40:214--225.
\url{https://doi.org/10.1016/0022-247X(72)90045-5}

\bibitem[Sharkey(1982)]{Sharkey1982}
Sharkey WW (1982) Cooperative games with large cores.
Int J Game Theory 11:175--182.
\url{https://doi.org/10.1007/BF01755727}

\end{thebibliography}

\end{document}
